\chapter{Thất Tình Đúng Mùa Thi}
\chapterintro{Có những nỗi buồn chọn thời điểm rất thiếu tinh thần hợp tác}{Thi cử đã căng, thất tình chen vào thì lịch học càng dễ tan.}

\begin{lucbatbox}{Lục bát: Buồn ngay tuần kiểm tra}
Đang mùa ôn tập gấp ghê\
Tim em lại chọn đường về đổ mưa\
Mở bài chưa được mấy dòng\
Đã nghe lòng nhớ một vòng chuyện xưa\
Đề thi chưa kịp làm vừa\
Nước buồn đã tới lưa thưa cả ngày
\end{lucbatbox}

\begin{tutuyetbox}{Thất ngôn tứ tuyệt: Vở cạnh bài nhạc}
Sách vở nằm đó rất ngoan\
Điện thoại bật nhạc miên man suốt chiều\
Nếu mà trái tim ít xiêu\
Có khi bài cũ đã nhiều phần xong
\end{tutuyetbox}

\begin{ghichu}
Thất tình đúng mùa thi là kiểu trùng lịch mà người trong cuộc không hề muốn, nhưng hậu quả thì phòng học phải gánh chung.
\end{ghichu}

\ngancach

\begin{lucbatbox}{Lục bát: Muốn học mà không yên}
Ngồi vào em cũng muốn làm\
Mà câu đầu đọc đã tràn bóng ai\
Bạn bè thương, bảo từ từ\
Thi xong rồi hẵng buồn dư cũng vừa\
Nhưng tim đâu biết nghe chừa\
Nên trang giấy cứ lưa thưa mấy hàng
\end{lucbatbox}

\begin{tutuyetbox}{Thất ngôn tứ tuyệt: Qua mùa rồi học}
Buồn thì vẫn cứ được buồn\
Nhưng xin đừng để luôn cuốn bài theo\
Thi qua, tình lắng một chiều\
Ngẩng lên còn phải thương điều tương lai
\end{tutuyetbox}